\documentclass[11pt]{article}
\usepackage[letterpaper,margin=1in]{geometry}
\usepackage{amsmath,amssymb,amsthm,lmodern,microtype}
\usepackage[T1]{fontenc}
\usepackage{xcolor,framed}
\usepackage[colorlinks=true,linkcolor=blue!60!black,urlcolor=blue!60!black]{hyperref}
\definecolor{shadecolor}{RGB}{240,244,250}
\newcommand{\R}{\mathbb R}
\newcommand{\T}{^{\mathsf T}}
\newcommand{\ones}{\mathbf 1}
\newcommand{\answer}[1]{\begin{shaded}\noindent\textbf{Result.} #1\end{shaded}}
\setlength{\parindent}{0pt}
\setlength{\parskip}{0.5em}
\hypersetup{pdftitle={Homework 4 - Worked Practice Solutions},pdfauthor={Optimization companion notes}}
\begin{document}
\begin{center}
{\Large\bfseries Optimization --- Homework 4}\\[5pt]
{\large Worked practice solutions}\\[5pt]
September 27, 2026
\end{center}
\textit{Scope.} Problems 1--4 and optional Problem 5 from the supplied handout.
Vector inequalities are componentwise; $\ones$ is an all-ones vector of the required dimension. These are companion explanations for practice.

\section*{Problem 1. A nonnegative invariant vector}
Let $P\in\R^{n\times n}$ be row stochastic: $P\ge0$ and $P\ones=\ones$.
We seek a nonzero $\lambda\ge0$ with $(P\T-I)\lambda=0$.

\begin{proof}
Consider the auxiliary primal--dual pair
\[
\begin{array}{ll@{\qquad\qquad}ll}
\text{maximize}&\ones\T\lambda&\text{minimize}&0\\
\text{subject to}&(P\T-I)\lambda=0&\text{subject to}&(P-I)y\ge\ones,\\
&\lambda\ge0&&y\in\R^n.
\end{array}
\]
The dual follows from the standard pair
$\max\{d\T z:Bz=h,z\ge0\}$ and
$\min\{h\T y:B\T y\ge d,y\text{ free}\}$.

The dual is infeasible. Indeed, for any $y$, choose $k$ such that
$y_k=\max_j y_j$. Nonnegativity and unit row sums give
\[
(Py)_k=\sum_j p_{kj}y_j
\le\sum_j p_{kj}y_k=y_k,
\]
so $((P-I)y)_k\le0$, contradicting the required lower bound $1$.

The primal is feasible at $\lambda=0$. If it had no nonzero feasible point,
its optimum would be the finite value $0$; strong LP duality would then
give a feasible dual optimum. This contradicts dual infeasibility.
Therefore a nonzero feasible $\lambda$ exists.

Since $\lambda\ge0$ and $\lambda\ne0$, its sum is positive. Hence
\[
\pi=\frac{\lambda}{\ones\T\lambda}
\quad\Longrightarrow\quad
P\T\pi=\pi,\qquad\pi\ge0,\qquad\ones\T\pi=1.
\]
Thus $\pi$ is an invariant probability distribution.
\end{proof}

\answer{Every finite row-stochastic matrix has a nonnegative left eigenvector for eigenvalue $1$, which can be normalized to a probability vector.}
\textbf{Interpretation.} $P\ones=\ones$ already gives a right eigenvector;
the proof supplies the nonnegative \emph{left} eigenvector needed for a
column-vector probability distribution. It does not assert uniqueness
or convergence of the Markov chain.

\newpage
\section*{Problem 2. A piecewise-linear convex objective}
Write
\[
f(x)=\max_{1\le i\le m}(a_i\T x-b_i),\qquad
v=\inf_{x\in\R^n}f(x)\in\R,
\]
where $m\ge1$, $A\in\R^{m\times n}$ has rows $a_i\T$, and $b\in\R^m$.

\subsection*{(a) Every feasible weight vector gives a lower bound}
\begin{proof}
Let $p\ge0$, $\ones\T p=1$, and $A\T p=0$. A weighted average is no
larger than the largest entry, so for every $x$,
\[
f(x)\ge\sum_i p_i(a_i\T x-b_i)
=x\T A\T p-p\T b=-p\T b.
\]
Taking the infimum over $x$ proves $-p\T b\le v$.
\end{proof}

\subsection*{(b) The best lower bound is exact}
\begin{proof}
Introduce a free scalar $t$ to obtain the equivalent epigraph LP
\[
\begin{aligned}
\min_{x\in\R^n,\ t\in\R}\quad&t\\
\text{s.t.}\quad&Ax-t\ones\le b.
\end{aligned}
\]
For fixed $x$, the smallest feasible $t$ is exactly $f(x)$.
For the inequality $Ax-t\ones-b\le0$, use multipliers $p\ge0$:
\[
\begin{aligned}
L(x,t,p)
&=t+p\T(Ax-t\ones-b)\\
&=(A\T p)\T x+(1-\ones\T p)t-p\T b.
\end{aligned}
\]
Because $x$ and $t$ are unrestricted, the infimum of $L$ over them is
finite precisely when $A\T p=0$ and $\ones\T p=1$. It is then $-p\T b$.
The dual is therefore
\[
\boxed{\max_p\ -p\T b\quad\text{s.t.}\quad
A\T p=0,\quad\ones\T p=1,\quad p\ge0.}
\]
The epigraph LP is feasible, for example at
$x=0$, $t=\max_i(-b_i)$, and its objective is bounded below by the
assumed finite value $v$. A feasible LP with a finite optimal value
attains it. Strong duality gives an attained dual optimum equal to $v$.
\end{proof}

\textbf{Connection to Homework 3.} The dual constructs a certificate by
combining constraint rows with nonnegative weights. Here the weights sum
to one because the coefficient of the free epigraph variable must vanish.
At an optimal pair, complementary slackness gives
\[
p_i^*\big[t^*-(a_i\T x^*-b_i)\big]=0.
\]
Only affine pieces that attain the maximum at $x^*$ can receive positive
weight. Their slopes balance: $\sum_i p_i^*a_i=0$.

\newpage
\section*{Problem 3. Symmetry produces an optimality certificate}
The primal and dual are
\[
\begin{array}{ll@{\qquad\qquad}ll}
\text{minimize}&c\T x&\text{maximize}&c\T y\\
\text{subject to}&Ax\ge c&\text{subject to}&A\T y\le c,\\
&x\ge0&&y\ge0.
\end{array}
\]
\begin{proof}
Suppose $Ax^*=c$ and $x^*\ge0$. Then $x^*$ is primal feasible.
Since $A=A\T$, the choice $y^*=x^*$ also satisfies
$A\T y^*=c\le c$ and $y^*\ge0$, so it is dual feasible.
Their objective values coincide: $c\T x^*=c\T y^*$.
Weak duality therefore certifies that both are optimal.
\end{proof}

\textbf{The same certificate without dual notation.} For any feasible $x$,
\[
c\T x=(Ax^*)\T x=(x^*)\T Ax
\ge(x^*)\T c=c\T x^*.
\]
The middle equality uses symmetry; the inequality uses $x^*\ge0$ and
$Ax\ge c$. No positive-semidefinite assumption on $A$ is needed.

\section*{Problem 4. Robust right-hand-side uncertainty}
The displayed uncertainty set is
\[
\mathcal U=\{b\in\R^m:0\le b_i\le1,\ \sum_i b_i\le1\}.
\]
\textit{Notation note.} Although the handout calls this a unit box,
the displayed set is the simplex $\operatorname{conv}\{0,e_1,\ldots,e_m\}$.
We use the displayed set; the claimed answer is also valid for $[0,1]^m$.

\begin{proof}
Suppose $Ax\ge b$ for every $b\in\mathcal U$. For each $i$, the coordinate
vector $e_i$ belongs to $\mathcal U$. Applying the robust constraint to
$b=e_i$ yields $(Ax)_i\ge1$. Thus $Ax\ge\ones$.

Conversely, if $Ax\ge\ones$, then every $b\in\mathcal U$ satisfies
$b\le\ones$, and hence $Ax\ge\ones\ge b$.
The feasible sets coincide and the objective is unchanged, proving
equivalence to $\min c\T x$ subject to $Ax\ge\ones$.
\end{proof}

\answer{The robust right-hand side is the vector of coordinatewise maxima:
$\widehat b_i=\max_{b\in\mathcal U}b_i=1$.}
For $m>1$, $\ones\notin\mathcal U$. This causes no problem: different
rows can have different worst-case scenarios. There is no sign restriction
on $x$ in this problem unless separately imposed.

\newpage
\section*{Problem 5 (optional). Finite robust counterparts}
\textbf{Reading the handout.} We use $Ax\le b$ with
$A\in\R^{m\times n}$ and $x\in\R^n$, consistent with its statement that
there are $m$ constraints. The printed $A\T x$ is dimensionally inconsistent
with those dimensions. Equivalently, one can keep $A\T x$ and take
$A\in\R^{n\times m}$. We call the last uncertainty set $\mathcal U_3$,
correcting its repeated $\mathcal U_2$ label.

Let
\[
\mathcal U_2=\{a\in\R^n:D_2a=d_2,\ a\ge0\},\qquad
\mathcal U_3=\{a\in\R^n:D_3a\ge d_3,\ a\le0\},
\]
where $D_j\in\R^{r_j\times n}$ and $d_j\in\R^{r_j}$.
First assume all uncertainty sets are nonempty.

\subsection*{1. Uncertain right-hand side}
Since $\mathcal U_1\subseteq\R^m$ is a nonempty polytope, compute the
finite constants
\[
\underline b_i=\min_{b\in\mathcal U_1}b_i,\qquad i=1,\ldots,m.
\]
Then $Ax\le b$ for every $b\in\mathcal U_1$ is equivalent to
$Ax\le\underline b$. If $\mathcal U_1=\{b:Gb\le h\}$, each
$\underline b_i$ is the value of the LP $\min\{e_i\T b:Gb\le h\}$.
If vertices $v^1,\ldots,v^K$ are given instead, use
$\underline b_i=\min_k v_i^k$.

\subsection*{2. Coefficients in $\mathcal U_2$}
For fixed $x$, the worst-case left-hand side is
\[
\begin{aligned}
s_2(x)&=\sup\{x\T a:D_2a=d_2,\ a\ge0\}\\
&=\min\{d_2\T y_2:D_2\T y_2\ge x,\ y_2\text{ free}\},
\end{aligned}
\]
whenever the supremum is finite, by LP strong duality.
Thus the robust inequality is equivalent to the existence of $y_2$ with
\[
\boxed{D_2\T y_2\ge x,\qquad d_2\T y_2\le b_2,\qquad y_2\in\R^{r_2}.}
\]
\begin{proof}[Why this is exact]
If such $y_2$ exists, every $a\in\mathcal U_2$ satisfies
$x\T a\le y_2\T D_2a=y_2\T d_2\le b_2$.
Conversely, if the robust inequality holds, the inner maximization is
nonempty and bounded above by $b_2$. Its dual therefore has an optimizer
that satisfies the displayed conditions.
\end{proof}

\newpage
\subsection*{3. Coefficients in $\mathcal U_3$}
Set $z=-a\ge0$. Then $D_3a\ge d_3$ becomes
$D_3z\le-d_3$, and $x\T a=(-x)\T z$. The worst-case LP and its dual are
\[
\begin{aligned}
s_3(x)&=\sup\{(-x)\T z:D_3z\le-d_3,\ z\ge0\}\\
&=\min\{-d_3\T\mu_3:D_3\T\mu_3\ge-x,\ \mu_3\ge0\},
\end{aligned}
\]
when the supremum is finite. Consequently,
\[
\boxed{D_3\T\mu_3\ge-x,\qquad-d_3\T\mu_3\le b_3,\qquad\mu_3\ge0.}
\]
\begin{proof}[Direct sign check]
For any $a\in\mathcal U_3$ and such a $\mu_3$, multiplying
$D_3\T\mu_3\ge-x$ by $a\le0$ reverses the inequality, giving
$a\T D_3\T\mu_3\le-a\T x$.
Also $\mu_3\T D_3a\ge\mu_3\T d_3$. Therefore
$x\T a\le-\mu_3\T D_3a\le-\mu_3\T d_3\le b_3$.
Necessity follows from strong duality as in the preceding case.
\end{proof}

\subsection*{The resulting finite LP}
For nonempty uncertainty sets, the complete formulation is
\[
\boxed{\begin{aligned}
\min_{x,y_2,\mu_3}\quad&c\T x\\
\text{s.t.}\quad&Ax\le\underline b,\\
&D_2\T y_2\ge x,\qquad d_2\T y_2\le b_2,\\
&D_3\T\mu_3\ge-x,\qquad-d_3\T\mu_3\le b_3,\\
&\mu_3\ge0,\qquad x\in\R^n,\quad y_2\in\R^{r_2}.
\end{aligned}}
\]
It has finitely many variables and linear constraints. The $m$ constants
$\underline b_i$ are precomputed by finite LPs.

\subsection*{Empty and unbounded uncertainty sets}
The handout does not explicitly assume nonemptiness. Check each set
by an LP feasibility problem first. If a set is empty, its universal
constraint is vacuous, so omit the corresponding block above.

$\mathcal U_2$ and $\mathcal U_3$ need not be bounded. For a given $x$,
an unbounded worst-case objective makes its robust constraint impossible
for finite $b_j$; the associated dual block then has no feasible
certificate. An unbounded \emph{set} need not have an unbounded objective
in every direction, so bounded uncertainty is not required for the
reformulation.

\vfill
\small\textit{Source:} supplied \texttt{hw04.pdf}, pp.\ 1--3.
Proofs and explanations are derived here; see the companion study guide
for geometric intuition, worked examples, self-tests, and a reference sheet.
\end{document}
